\SourceLeaf{250}%
\noindent entirely into triangular pyramids that are pairwise symmetric in the two solids.

II. A sphere can be constructed about each triangular pyramid, whose 4 radii divide it by means of planes drawn through each pair into 4 equilateral pyramids. (Allowing for the possibility that there may also be negative pyramids.)

III. If this construction is carried out for each pair of triangular pyramids mentioned under I., it gives equilateral pyramids throughout, again symmetric in pairs.

IV. Every equilateral pyramid can be divided into \textls{three} triangular pyramids by means of a perpendicular from the apex to the base and planes drawn through it. The base of each is an isosceles triangle, and the edge issuing from the apex of this isosceles triangle is perpendicular to the base.

V. Two pyramids of the kind just described that are situated correspondingly in the two polyhedra are both symmetric and congruent, because each can be divided by a plane into two halves symmetric to one another.

VI. Thus the two polyhedra can be regarded as the (algebraic) sums of congruent pyramids throughout, so they have equal volumes. Q.E.D.

I must add that the proof is only partly my own. In thinking over your remark, I had got far enough to see parts III---V clearly, but at first I found it difficult to reduce every triangular pyramid to equilateral ones. I might have let the matter drop had it not come up by chance in a conversation with Dr.\ \textsc{Stegmann}, who had the fortunate idea of using the sphere around the triangular [pyramid]. This really provided the key.

As far as I can see, there is no fallacy in the foregoing. I should be much obliged if you would examine the proof closely and write to me on some occasion with your judgment of it. \mbox{.\ .\ .\ .\ .\ .\ .}\}
